\documentclass[aspectratio=169,10pt]{beamer}
\usetheme{metropolis}
\usepackage{booktabs,amsmath,amssymb}
\title{Q15: Context, Retrieval, and Memory}
\subtitle{Windows, RAG, and the volatility of context}
\author{DLS 2026 --- Kukla}
\date{October 2026}
\begin{document}
\maketitle

\begin{frame}{The question}
\begin{block}{Student question (\#13, \#15)}
How do context windows, retrieval, and conversation memory work, and how do we avoid hallucination without
deleting context?
\end{block}
\end{frame}
\begin{frame}{The context window + its cost}
\begin{itemize}
\item All the model ``knows now'' must fit in $n$ tokens.
\item Attention compute is $O(n^2)$; \textbf{KV cache} grows linearly in $n$.
\item ``Just make it bigger'' is \emph{not} free.
\end{itemize}
\end{frame}
\begin{frame}{Retrieval-augmented generation}
embed docs $\to$ index $\to$ embed query $\to$ top-$k$ $\to$ put in context.
\begin{itemize}
\item Draw on a corpus far larger than any window.
\item Pay attention cost only on retrieved chunks.
\item \textbf{Cite grounded sources} $\Rightarrow$ less hallucination.
\end{itemize}
\end{frame}
\begin{frame}{Long context vs.\ retrieval}
\begin{itemize}
\item \textbf{Long context}: coherent whole, reason jointly (contract, code file).
\item \textbf{Retrieval}: vast, sparse, fresh, or must-attribute knowledge.
\item Common pattern: retrieve to narrow, then long-window reason.
\end{itemize}
\end{frame}
\begin{frame}{Why context is not memory}
\begin{itemize}
\item Context = \textbf{volatile} working memory; gone when chat ends / scrolls out.
\item Never written back to weights (Q14's long-term store).
\item Agent memory: persist to a store, summarize, \emph{re-inject} selectively.
\end{itemize}
\alert{``Don't delete context'' = what to persist/summarize/retrieve, not infinite windows.}
\end{frame}

\end{document}
